\chapter{Protocol Design and Implementation}

This chapter presents the full version-by-version protocol evolution. The
implementation strategy is incremental. Each new version keeps previous
capabilities and introduces one or more carefully selected features. This
approach improves correctness and makes debugging practical.

\section{Global Design Principles}
The same design principles are followed across all versions.

\begin{enumerate}
    \item \textbf{Single responsibility per stage:} each version focuses on a
    specific protocol concept.
    \item \textbf{Explicit state:} sender and receiver maintain clearly named
    variables for sequence numbers, timeouts, and window state.
    \item \textbf{Deterministic transitions:} packet receipt leads to a defined
    state update path.
    \item \textbf{Traceable behavior:} console output helps validate the
    protocol timeline in each test run.
\end{enumerate}

\clearpage
\section{Version Family A: Bit and Byte Transfer (v0--v2)}
The first family introduces communication mechanics at the lowest useful level.

\subsection{Version 0: Single-Bit UDP Demo}
In level 0, the sender sends one bit to the receiver using a UDP datagram.
There is no framing beyond the single value. This version validates:

\begin{enumerate}
    \item socket creation,
    \item destination addressing,
    \item datagram transmission,
    \item datagram reception.
\end{enumerate}

This step is intentionally minimal and prepares the path for structured packet
design.

\vspace*{1.5\baselineskip}
\begin{figure}[h!]
\centering
\resizebox{\textwidth}{!}{%
\begin{tikzpicture}[
    node distance=1.4cm,
    sender/.style={draw, rounded corners=3pt, thick, minimum width=4.1cm, minimum height=1.0cm, align=center, fill=blue!10, draw=blue!60!black},
    receiver/.style={draw, rounded corners=3pt, thick, minimum width=4.1cm, minimum height=1.0cm, align=center, fill=teal!10, draw=teal!60!black},
    channel/.style={draw, rounded corners=3pt, thick, minimum width=3.8cm, minimum height=0.9cm, align=center, fill=orange!10, draw=orange!70!black},
    arr/.style={-{Latex[length=3mm]}, thick, draw=black!80},
    flow/.style={-{Latex[length=3mm]}, very thick, draw=red!70!black}
]
\node[sender] (sapp) at (0,1.8) {Sender Application\\(bit input: 0 or 1)};
\node[sender] (sudp) at (0,0) {Sender UDP Socket\\(sendto)};

\node[channel] (net) at (6,0) {Unreliable UDP Network};

\node[receiver] (rudp) at (12,0) {Receiver UDP Socket\\(recvfrom)};
\node[receiver] (rapp) at (12,1.8) {Receiver Application\\(bit output)};

\draw[arr] (sapp) -- (sudp);
\draw[flow] (sudp.east) -- (net.west);
\draw[flow] (net.east) -- (rudp.west);
\draw[arr] (rudp) -- (rapp);

\node[align=center] at (6,-1.35) {No framing, no ACK/NAK, no retransmission\\(only basic UDP send and receive path)};
\end{tikzpicture}%
}
\caption{Version 0 architecture: single-bit UDP transmission from sender to receiver.
}
\label{fig:v0-architecture}
\end{figure}

\clearpage
\subsection{Version 1: Character as 8 Independent Bits}
In level 1, one character is encoded into 8 bits and sent as 8 UDP datagrams.
The receiver reconstructs the byte by shifting and OR operations. Timing is
regulated using small sleep intervals so packet order remains easier to track
in local runs.

Educational value:
\begin{enumerate}
    \item understanding of binary representation,
    \item significance of bit order,
    \item reconstruction logic at receiver side.
\end{enumerate}

\vspace{1.25\baselineskip}
\begin{figure}[h!]
\centering
\resizebox{\textwidth}{!}{%
\begin{tikzpicture}[
    node distance=1.3cm,
    sender/.style={draw, rounded corners=3pt, thick, minimum width=4.4cm, minimum height=1.0cm, align=center, fill=blue!10, draw=blue!60!black},
    receiver/.style={draw, rounded corners=3pt, thick, minimum width=4.6cm, minimum height=1.0cm, align=center, fill=teal!10, draw=teal!60!black},
    channel/.style={draw, rounded corners=3pt, thick, minimum width=4.0cm, minimum height=0.9cm, align=center, fill=orange!10, draw=orange!70!black},
    arr/.style={-{Latex[length=3mm]}, thick, draw=black!80},
    flow/.style={-{Latex[length=3mm]}, very thick, draw=red!70!black}
]
\node[sender] (sapp) at (0,2.0) {Sender Application\\(input: one character)};
\node[sender] (senc) at (0,0.4) {Bit Encoder\\(split byte into b7...b0)};
\node[sender] (sudp) at (0,-1.2) {Sender UDP Socket\\(8 sendto calls)};

\node[channel] (net) at (6,-1.2) {Unreliable UDP Network};

\node[receiver] (rudp) at (12,-1.2) {Receiver UDP Socket\\(8 recvfrom calls)};
\node[receiver] (rdec) at (12,0.4) {Byte Reconstructor\\(shift + OR accumulation)};
\node[receiver] (rapp) at (12,2.0) {Receiver Application\\(restored character)};

\draw[arr] (sapp) -- (senc);
\draw[arr] (senc) -- (sudp);
\draw[flow] (sudp.east) -- (net.west);
\draw[flow] (net.east) -- (rudp.west);
\draw[arr] (rudp) -- (rdec);
\draw[arr] (rdec) -- (rapp);

\node[align=center] at (6,-2.45) {Datagram stream: bit\_7, bit\_6, ..., bit\_0\\Short sleep intervals help preserve local ordering during experiments};
\end{tikzpicture}%
}
\caption{Version 1 architecture: one character transmitted as eight UDP bit-datagrams and reconstructed at receiver side.
}
\label{fig:v1-architecture}
\end{figure}

\clearpage
\subsection{Version 2: Full String by Bit Stream}
Level 2 extends the bit-stream approach to full strings. A null character is
sent at the end to mark message completion. At this stage, communication works
for complete user input, but no error detection or reliability control exists.

\vspace{0.75\baselineskip}
\begin{figure}[h!]
\centering
\resizebox{\textwidth}{!}{%
\begin{tikzpicture}[
    node distance=1.3cm,
    sender/.style={draw, rounded corners=3pt, thick, minimum width=4.9cm, minimum height=1.0cm, align=center, fill=blue!10, draw=blue!60!black},
    receiver/.style={draw, rounded corners=3pt, thick, minimum width=4.9cm, minimum height=1.0cm, align=center, fill=teal!10, draw=teal!60!black},
    channel/.style={draw, rounded corners=3pt, thick, minimum width=4.0cm, minimum height=0.9cm, align=center, fill=orange!10, draw=orange!70!black},
    arr/.style={-{Latex[length=3mm]}, thick, draw=black!80},
    flow/.style={-{Latex[length=3mm]}, very thick, draw=red!70!black}
]
\node[sender] (sapp) at (0,2.0) {Sender Application\\(input: full string)};
\node[sender] (sstream) at (0,0.4) {Bit Stream Builder\\(all characters to bit stream + NULL terminator)};
\node[sender] (sudp) at (0,-1.2) {Sender UDP Socket\\(multiple sendto calls)};

\node[channel] (net) at (6,-1.2) {Unreliable UDP Network};

\node[receiver] (rudp) at (12,-1.2) {Receiver UDP Socket\\(multiple recvfrom calls)};
\node[receiver] (rparse) at (12,0.4) {Stream Parser\\(rebuild bytes until NULL detected)};
\node[receiver] (rapp) at (12,2.0) {Receiver Application\\(complete message output)};

\draw[arr] (sapp) -- (sstream);
\draw[arr] (sstream) -- (sudp);
\draw[flow] (sudp.east) -- (net.west);
\draw[flow] (net.east) -- (rudp.west);
\draw[arr] (rudp) -- (rparse);
\draw[arr] (rparse) -- (rapp);

\node[align=center] at (6,-2.45) {Transfer unit remains bit-level datagrams; end-of-message is marked by NULL character\\No integrity check and no reliability loop at this stage};
\end{tikzpicture}%
}
\caption{Version 2 architecture: full string transmitted as a bit stream with NULL-terminated completion.
}
\label{fig:v2-architecture}
\end{figure}

\clearpage
\section{Version Family B: Error Detection and Stop-and-Wait (v3--v7)}
Family B introduces reliability fundamentals.
\subsection{Version 3: Parity Detection}
For each byte, sender computes parity and transmits it with payload bits.
Receiver recomputes parity and compares values.

Outcome:
\begin{enumerate}
    \item corruption can be detected,
    \item corruption cannot yet be corrected.
\end{enumerate}

Parity is simple and fast but weak against many multi-bit error patterns.

\vspace{1.5\baselineskip}
\begin{figure}[h!]
\centering
\resizebox{\textwidth}{!}{%
\begin{tikzpicture}[
    node distance=1.3cm,
    sender/.style={draw, rounded corners=3pt, thick, minimum width=4.7cm, minimum height=1.0cm, align=center, fill=blue!10, draw=blue!60!black},
    receiver/.style={draw, rounded corners=3pt, thick, minimum width=4.9cm, minimum height=1.0cm, align=center, fill=teal!10, draw=teal!60!black},
    channel/.style={draw, rounded corners=3pt, thick, minimum width=4.1cm, minimum height=0.9cm, align=center, fill=orange!10, draw=orange!70!black},
    arr/.style={-{Latex[length=3mm]}, thick, draw=black!80},
    flow/.style={-{Latex[length=3mm]}, very thick, draw=red!70!black}
]
\node[sender] (sapp) at (0,2.0) {Sender Application\\(payload byte input)};
\node[sender] (spar) at (0,0.4) {Parity Generator\\(compute parity bit)};
\node[sender] (sudp) at (0,-1.2) {Sender UDP Socket\\(send payload + parity)};

\node[channel] (net) at (6,-1.2) {Unreliable UDP Network};

\node[receiver] (rudp) at (12,-1.2) {Receiver UDP Socket\\(receive payload + parity)};
\node[receiver] (rchk) at (12,0.4) {Parity Checker\\(recompute and compare)};
\node[receiver] (rout) at (12,2.0) {Receiver Output\\(accept or flag corruption)};

\draw[arr] (sapp) -- (spar);
\draw[arr] (spar) -- (sudp);
\draw[flow] (sudp.east) -- (net.west);
\draw[flow] (net.east) -- (rudp.west);
\draw[arr] (rudp) -- (rchk);
\draw[arr] (rchk) -- (rout);

\node[align=center] at (6,-2.45) {Parity mismatch indicates corruption can be detected\\No correction and no retransmission in this version};
\end{tikzpicture}%
}
\caption{Version 3 architecture: parity-based corruption detection over UDP payload transfer.
}
\label{fig:v3-architecture}
\end{figure}

\clearpage
\subsection{Version 4: Sequence Numbers}
Level 4 attaches sequence numbers to per-character transfer. This improves
observability and creates the basis for reliable acknowledgement logic.

\vspace{1.5\baselineskip}
\begin{figure}[h!]
\centering
\resizebox{\textwidth}{!}{%
\begin{tikzpicture}[
    node distance=1.3cm,
    sender/.style={draw, rounded corners=3pt, thick, minimum width=5.0cm, minimum height=1.0cm, align=center, fill=blue!10, draw=blue!60!black},
    receiver/.style={draw, rounded corners=3pt, thick, minimum width=5.0cm, minimum height=1.0cm, align=center, fill=teal!10, draw=teal!60!black},
    channel/.style={draw, rounded corners=3pt, thick, minimum width=4.1cm, minimum height=0.9cm, align=center, fill=orange!10, draw=orange!70!black},
    arr/.style={-{Latex[length=3mm]}, thick, draw=black!80},
    flow/.style={-{Latex[length=3mm]}, very thick, draw=red!70!black}
]
\node[sender] (sapp) at (0,2.0) {Sender Application\\(character stream)};
\node[sender] (spkt) at (0,0.4) {Sequencer\\(attach seq number to each unit)};
\node[sender] (sudp) at (0,-1.2) {Sender UDP Socket\\(send \{seq, payload\})};

\node[channel] (net) at (6,-1.2) {Unreliable UDP Network};

\node[receiver] (rudp) at (12,-1.2) {Receiver UDP Socket\\(receive \{seq, payload\})};
\node[receiver] (rord) at (12,0.4) {Sequence Observer\\(track arrival order by seq)};
\node[receiver] (rapp) at (12,2.0) {Receiver Output\\(ordered view and diagnostics)};

\draw[arr] (sapp) -- (spkt);
\draw[arr] (spkt) -- (sudp);
\draw[flow] (sudp.east) -- (net.west);
\draw[flow] (net.east) -- (rudp.west);
\draw[arr] (rudp) -- (rord);
\draw[arr] (rord) -- (rapp);

\node[align=center] at (6,-2.45) {Sequence numbers increase observability of packet order and duplicates\\This stage prepares state tracking needed for ACK-based reliability};
\end{tikzpicture}%
}
\caption{Version 4 architecture: sequence-numbered UDP units for order visibility and reliability groundwork.
}
\label{fig:v4-architecture}
\end{figure}

\clearpage
\subsection{Version 5: Stop-and-Wait ARQ with ACK/NAK}
In level 5, receiver sends ACK on success and NAK on parity failure. Sender
transmits one unit and waits for response before moving forward. This is the
first true reliability loop.

Strength:
\begin{enumerate}
    \item simple correctness argument,
    \item easy handling of duplicates,
    \item straightforward debugging.
\end{enumerate}

Limitation:
\begin{enumerate}
    \item low throughput because channel stays idle while waiting.
\end{enumerate}

\vspace{1.5\baselineskip}
\begin{figure}[h!]
\centering
\resizebox{\textwidth}{!}{%
\begin{tikzpicture}[
    node distance=1.3cm,
    sender/.style={draw, rounded corners=3pt, thick, minimum width=5.0cm, minimum height=1.0cm, align=center, fill=blue!10, draw=blue!60!black},
    receiver/.style={draw, rounded corners=3pt, thick, minimum width=5.0cm, minimum height=1.0cm, align=center, fill=teal!10, draw=teal!60!black},
    channel/.style={draw, rounded corners=3pt, thick, minimum width=4.1cm, minimum height=0.9cm, align=center, fill=orange!10, draw=orange!70!black},
    arr/.style={-{Latex[length=3mm]}, thick, draw=black!80},
    flow/.style={-{Latex[length=3mm]}, very thick, draw=red!70!black},
    ack/.style={-{Latex[length=3mm]}, very thick, draw=green!55!black}
]
\node[sender] (sapp) at (0,2.0) {Sender Logic\\(send one unit and wait)};
\node[sender] (swait) at (0,0.4) {Stop-and-Wait Controller\\(state: wait for ACK/NAK)};
\node[sender] (sudp) at (0,-1.2) {Sender UDP Socket\\(send DATA unit)};

\node[channel] (net) at (6,-1.2) {Unreliable UDP Network};

\node[receiver] (rudp) at (12,-1.2) {Receiver UDP Socket\\(receive DATA unit)};
\node[receiver] (rcheck) at (12,0.4) {Parity Validation\\(success or failure)};
\node[receiver] (rresp) at (12,2.0) {Response Generator\\(ACK on success, NAK on error)};

\draw[arr] (sapp) -- (swait);
\draw[arr] (swait) -- (sudp);
\draw[flow] (sudp.east) -- node[above, fill=white, inner sep=1pt]{DATA} (net.west);
\draw[flow] (net.east) -- (rudp.west);
\draw[arr] (rudp) -- (rcheck);
\draw[arr] (rcheck) -- (rresp);

\draw[ack] (rresp.west) -- ++(-0.25,0) |- node[pos=0.22, left, fill=white, inner sep=1pt]{ACK/NAK} (swait.east);

\node[align=center] at (6,-2.45) {First reliability loop: sender advances only after ACK\\On NAK, sender retransmits the same unit};
\end{tikzpicture}%
}
\caption{Version 5 architecture: stop-and-wait ARQ with ACK/NAK feedback path.
}
\label{fig:v5-architecture}
\end{figure}

\clearpage
\subsection{Version 6: Timeout and Retries}
Level 6 adds socket receive timeout and retry limit. If ACK does not arrive in
time, sender retransmits. Receiver tracks expected sequence and handles
duplicates conservatively.

This version is the first practically robust design under packet loss.

\vspace{1.5\baselineskip}
\begin{figure}[h!]
\centering
\resizebox{\textwidth}{!}{%
\begin{tikzpicture}[
    node distance=1.3cm,
    sender/.style={draw, rounded corners=3pt, thick, minimum width=5.2cm, minimum height=1.0cm, align=center, fill=blue!10, draw=blue!60!black},
    receiver/.style={draw, rounded corners=3pt, thick, minimum width=5.2cm, minimum height=1.0cm, align=center, fill=teal!10, draw=teal!60!black},
    channel/.style={draw, rounded corners=3pt, thick, minimum width=4.1cm, minimum height=0.9cm, align=center, fill=orange!10, draw=orange!70!black},
    arr/.style={-{Latex[length=3mm]}, thick, draw=black!80},
    flow/.style={-{Latex[length=2mm]}, very thick, draw=red!70!black},
    ack/.style={-{Latex[length=2mm]}, very thick, draw=green!55!black},
    retry/.style={-{Latex[length=3mm]}, very thick, draw=purple!70!black}
]
\node[sender] (sctrl) at (0,2.0) {Sender Reliability Controller\\(DATA send, timeout timer, retry limit)};
\node[sender] (sstate) at (0,0.4) {State Machine\\(send to wait ACK to timeout to retransmit)};
\node[sender] (sudp) at (0,-1.2) {Sender UDP Socket\\(DATA / retransmitted DATA)};

\node[channel] (net) at (6,-1.2) {Unreliable UDP Network};

\node[receiver] (rudp) at (12,-1.2) {Receiver UDP Socket\\(DATA receive)};
\node[receiver] (rseq) at (12,0.4) {Expected-Seq Tracker\\(accept new, identify duplicates)};
\node[receiver] (rack) at (12,2.0) {ACK Generator\\(ACK valid expected unit)};

\draw[arr] (sctrl) -- (sstate);
\draw[arr] (sstate) -- (sudp);
\draw[flow] (sudp.east) -- node[above, fill=white, inner sep=1pt]{\small DATA} ++(0.4,0) -- (net.west);
\draw[flow] (net.east) -- (rudp.west);
\draw[arr] (rudp) -- (rseq);
\draw[arr] (rseq) -- (rack);
\draw[ack] (rack.west) -- ++(-3.5,0) |- node[pos=0.24, left, fill=white, inner sep=1pt]{ACK} (sstate.east);

\node[align=center] at (6,-2.55) {Timeout prevents indefinite waiting; retries improve robustness under loss\\Receiver sequence tracking prevents duplicate delivery};
\end{tikzpicture}%
}
\caption{Version 6 architecture: timeout-driven retransmission with duplicate-safe receiver handling.
}
\label{fig:v6-architecture}
\end{figure}

\clearpage
\subsection{Version 7: Struct-Based Packetization}
Version 7 replaces loose fields with an explicit packet structure:
\begin{enumerate}
    \item sequence number,
    \item payload,
    \item parity,
    \item end flag.
\end{enumerate}

This change improves code readability and prepares the architecture for
multi-field TCP-like packets in later levels.

\vspace{1.5\baselineskip}
\begin{figure}[h!]
\centering
\resizebox{\textwidth}{!}{%
\begin{tikzpicture}[
    node distance=1.35cm,
    sender/.style={draw, rounded corners=3pt, thick, minimum width=5.8cm, minimum height=1.1cm, align=center, fill=blue!10, draw=blue!60!black},
    receiver/.style={draw, rounded corners=3pt, thick, minimum width=5.8cm, minimum height=1.1cm, align=center, fill=teal!10, draw=teal!60!black},
    channel/.style={draw, rounded corners=3pt, thick, minimum width=4.6cm, minimum height=1.0cm, align=center, fill=orange!10, draw=orange!70!black},
    arr/.style={-{Latex[length=3mm]}, thick, draw=black!80},
    flow/.style={-{Latex[length=3mm]}, line width=1.1pt, draw=red!70!black},
    ack/.style={-{Latex[length=3mm]}, line width=1.1pt, draw=green!55!black}
]
\node[sender] (fields) at (-1.6,2.0) {Input Fields\\(sequence, payload, parity, end flag)};
\node[sender] (pack) at (-1.6,0.4) {Packet Struct Builder\\(pack all fields into one frame)};
\node[sender] (send) at (-1.6,-1.2) {Serialized Packet Sender\\(single structured packet)};

\node[channel] (net) at (7.6,-1.2) {UDP Channel Carrying Structured Frame};

\node[receiver] (recv) at (16.8,-1.2) {Packet Receiver\\(one frame arrives intact)};
\node[receiver] (parse) at (16.8,0.4) {Struct Parser\\(extract seq, payload, parity, end)};
\node[receiver] (validate) at (16.8,2.0) {Validation + Delivery\\(integrity and order checks)};

\draw[arr] (fields) -- (pack);
\draw[arr] (pack) -- (send);
\draw[flow] (send.east) -- node[pos=0.5, above, fill=white, inner sep=1pt]{packet struct} (net.west);
\draw[flow] (net.east) -- node[pos=0.5, above, fill=white, inner sep=1pt]{structured frame} (recv.west);
\draw[arr] (recv) -- (parse);
\draw[arr] (parse) -- (validate);
\draw[ack] (validate.west) -- ++(-1.8,0) |- node[pos=0.30, left, fill=white, inner sep=1pt]{status / ACK} (pack.east);

\node[align=center] at (7.6,-2.75) {Version 7 unifies metadata and payload into one explicit structure\\Cleaner parsing reduces mismatch between sender and receiver logic};
\end{tikzpicture}%
}
\caption{Version 7 architecture: struct-based packetization with unified sender and receiver parsing flow.
}
\label{fig:v7-architecture}
\end{figure}

\clearpage
\section{Version Family C: TCP-like Framing and Pipelining (v8--v10)}
This family introduces transport semantics closer to TCP.

\subsection{Version 8: Handshake, Flags, and Connection Lifecycle}
Level 8 includes:
\begin{enumerate}
    \item SYN for start,
    \item ACK for acknowledgement,
    \item DATA for payload,
    \item FIN for graceful close.
\end{enumerate}

The sender and receiver perform a simplified three-way handshake and teardown.
Checksum is also introduced at this level for stronger integrity checks.

\vspace{1.5\baselineskip}
\begin{figure}[h!]
\centering
\resizebox{\textwidth}{!}{%
\begin{tikzpicture}[
    font=\normalsize,
    sender/.style={draw, rounded corners=3pt, thick, minimum width=5.6cm, minimum height=1.1cm, align=center, fill=blue!10, draw=blue!60!black},
    receiver/.style={draw, rounded corners=3pt, thick, minimum width=5.6cm, minimum height=1.1cm, align=center, fill=teal!10, draw=teal!60!black},
    channel/.style={draw, rounded corners=3pt, thick, minimum width=4.0cm, minimum height=1.0cm, align=center, fill=orange!10, draw=orange!70!black},
    ctrl/.style={-{Latex[length=3.6mm]}, line width=1.0pt, draw=black!80},
    syn/.style={-{Latex[length=3.6mm]}, line width=1.25pt, draw=blue!70!black},
    ack/.style={-{Latex[length=3.6mm]}, line width=1.25pt, draw=green!55!black},
    data/.style={-{Latex[length=3.6mm]}, line width=1.25pt, draw=red!70!black},
    fin/.style={-{Latex[length=3.6mm]}, line width=1.25pt, draw=purple!70!black}
]
\node[sender] (shs) at (-2.4,2.1) {Sender Connection State\\(CLOSED to SYN-SENT to ESTABLISHED to FIN-WAIT)};
\node[sender] (sfrm) at (-2.4,0.5) {Sender Frame Builder\\(SYN/ACK/DATA/FIN flags + checksum)};
\node[sender] (snet) at (-2.4,-1.1) {Sender Transport Output\\(flagged frames to network)};

\node[channel] (net) at (8.1,-1.1) {Transport Channel for Flagged Frames};

\node[receiver] (rnet) at (18.4,-1.1) {Receiver Transport Input\\(accept flagged frames)};
\node[receiver] (rparse) at (18.4,0.5) {Receiver Parser + Checksum Verify\\(drop corrupted, accept valid)};
\node[receiver] (rstate) at (18.4,2.1) {Receiver Connection State\\(LISTEN to ESTABLISHED to CLOSE-WAIT)};

\draw[ctrl] (shs) -- (sfrm);
\draw[ctrl] (sfrm) -- (snet);
\draw[ctrl] (rnet) -- (rparse);
\draw[ctrl] (rparse) -- (rstate);

\draw[syn] (snet.east) -- node[pos=0.5, above, fill=white, inner sep=1pt]{SYN} (net.west);
\draw[syn] (net.east) -- node[pos=0.5, above, fill=white, inner sep=1pt]{SYN} (rnet.west);

\draw[ack] (rstate.north west) -- ++(-4.0,0) |- node[pos=0.34, left, fill=white, inner sep=1pt]{SYN ACK / ACK} (sfrm.north east);

\draw[data] (snet.south east) -- (2.6,-2.0) -- node[pos=0.50, below, fill=white, inner sep=1pt]{DATA + checksum} (13.4,-2.0) -- (rnet.south west);

\draw[fin] (sfrm.east) -- ++(2.4,0) |- node[pos=0.62, above, fill=white, inner sep=1pt]{FIN teardown} (rstate.west);

\node[align=center] at (8.1,-3.35) {Version 8 adds handshake lifecycle (SYN, ACK, FIN) and per-frame checksum verification\\This is the first step toward TCP-like connection semantics in the project};
\end{tikzpicture}%
}
\caption{Version 8 architecture: handshake flags, lifecycle control, and checksum-protected data flow.
}
\label{fig:v8-architecture}
\end{figure}

\clearpage
\subsection{Version 9: Sliding Window (Go-Back-N Style)}
Version 9 increases efficiency with a sender window (typical size 4). Sender
may transmit multiple packets before waiting for ACKs. On timeout, packets from
base are retransmitted. Receiver sends cumulative-like ACK behavior based on
expected sequence.

This level marks a major throughput jump because waiting time is amortized.

\vspace{1.5\baselineskip}
\begin{figure}[h!]
\centering
\resizebox{\textwidth}{!}{%
\begin{tikzpicture}[
    font=\normalsize,
    sender/.style={draw, rounded corners=3pt, thick, minimum width=5.8cm, minimum height=1.15cm, align=center, fill=blue!10, draw=blue!60!black},
    receiver/.style={draw, rounded corners=3pt, thick, minimum width=5.8cm, minimum height=1.15cm, align=center, fill=teal!10, draw=teal!60!black},
    channel/.style={draw, rounded corners=3pt, thick, minimum width=4.4cm, minimum height=0.95cm, align=center, fill=orange!10, draw=orange!70!black},
    ctrl/.style={-{Latex[length=3.4mm]}, line width=1.0pt, draw=black!80},
    data/.style={-{Latex[length=3.4mm]}, line width=1.2pt, draw=red!70!black},
    ack/.style={-{Latex[length=3.4mm]}, line width=1.2pt, draw=green!55!black},
    retry/.style={-{Latex[length=3.4mm]}, line width=1.2pt, draw=purple!70!black}
]
\node[sender] (swin) at (-2.2,2.0) {Sender Window Controller\\(window size = 4, base and next sequence)};
\node[sender] (squeue) at (-2.2,0.4) {Pipeline Send Queue\\(seq i, i+1, i+2, i+3 in flight)};
\node[sender] (stimer) at (-2.2,-1.2) {Timeout and Retransmit Logic\\(on timeout, resend from base)};

\node[channel] (net) at (8.0,-1.2) {Unreliable Network Channel};

\node[receiver] (rbuf) at (18.0,-1.2) {Receiver Input + Ordering Check\\(accept expected, discard out-of-window)};
\node[receiver] (rack) at (18.0,0.4) {Cumulative ACK Generator\\(ACK highest contiguous sequence)};
\node[receiver] (rapp) at (18.0,2.0) {In-order Delivery Layer\\(release contiguous data to application)};

\draw[ctrl] (swin) -- (squeue);
\draw[ctrl] (squeue) -- (stimer);
\draw[ctrl] (rbuf) -- (rack);
\draw[ctrl] (rack) -- (rapp);

\draw[data] (squeue.east) -- node[pos=0.5, above, fill=white, inner sep=1pt]{windowed DATA packets} (net.west);
\draw[data] (net.east) -- node[pos=0.5, above, fill=white, inner sep=1pt]{multiple in-flight frames} (rbuf.west);

\draw[ack] (rack.north west) -- ++(-4.0,0) |- node[pos=0.34, left, fill=white, inner sep=1pt]{cumulative ACK} (swin.north east);

\draw[retry] (stimer.south east) -- (2.0,-2.1) -- node[pos=0.48, below, fill=white, inner sep=1pt]{timeout retransmit from base} (13.0,-2.1) -- (rbuf.south west);

\node[align=center] at (8.0,-3.35) {Version 9 pipelines multiple packets before ACK wait, improving throughput\\Cumulative ACK plus timeout-from-base retransmission defines Go-Back-N behavior};
\end{tikzpicture}%
}
\caption{Version 9 architecture: sliding window pipelining with cumulative ACK and timeout-based retransmission.
}
\label{fig:v9-architecture}
\end{figure}

\clearpage
\subsection{Version 10: Chunk-Based Payloads}
Version 10 increases data per packet by introducing \texttt{CHUNK\_SIZE = 8}
bytes. Benefits:

\begin{enumerate}
    \item fewer packets for same message length,
    \item less per-packet header overhead,
    \item better throughput for moderate messages.
\end{enumerate}

Receiver logic starts supporting in-window buffering, enabling transition to
selective-repeat style delivery in the next family.

\vspace{1.5\baselineskip}
\begin{figure}[h!]
\centering
\resizebox{\textwidth}{!}{%
\begin{tikzpicture}[
    font=\normalsize,
    sender/.style={draw, rounded corners=3pt, thick, minimum width=5.8cm, minimum height=1.05cm, align=center, fill=blue!10, draw=blue!60!black},
    receiver/.style={draw, rounded corners=3pt, thick, minimum width=5.8cm, minimum height=1.05cm, align=center, fill=teal!10, draw=teal!60!black},
    channel/.style={draw, rounded corners=3pt, thick, minimum width=4.4cm, minimum height=0.95cm, align=center, fill=orange!10, draw=orange!70!black},
    ctrl/.style={-{Latex[length=3.4mm]}, line width=1.0pt, draw=black!80},
    data/.style={-{Latex[length=3.4mm]}, line width=1.2pt, draw=red!70!black},
    ack/.style={-{Latex[length=3.4mm]}, line width=1.2pt, draw=green!55!black},
    flow2/.style={-{Latex[length=3.4mm]}, line width=1.2pt, draw=purple!70!black}
]
\node[sender] (schunk) at (-3.0,2.6) {Sender Chunk Builder\\(split payload into 8-byte chunks)};
\node[sender] (spkt) at (-3.0,0.5) {Chunk Packetizer\\(seq + chunk + checksum)};
\node[sender] (sout) at (-3.0,-1.6) {Pipeline Output\\(fewer, larger packets for same message)};

\node[channel] (net) at (8.0,-1.6) {Transport Channel for Chunk Packets};

\node[receiver] (rin) at (19.0,-1.6) {Chunk Packet Receiver\\(validate and accept chunk packets)};
\node[receiver] (rbuf) at (19.0,0.5) {In-Window Buffer Support\\(store out-of-order chunks)};
\node[receiver] (rasm) at (19.0,2.6) {Chunk Reassembly Layer\\(reconstruct full message in order)};

\draw[ctrl] (schunk) -- (spkt);
\draw[ctrl] (spkt) -- (sout);
\draw[ctrl] (rin) -- (rbuf);
\draw[ctrl] (rbuf) -- (rasm);

\draw[data] (sout.east) -- node[pos=0.5, above, align=center, fill=white, inner sep=1pt]{chunk packets\\(8-byte payload units)} (net.west);
\draw[data] (net.east) -- node[pos=0.5, above, align=center, fill=white, inner sep=1pt]{reduced packet count\\higher payload efficiency} (rin.west);

\draw[ack] (rbuf.north west) -- ++(-4.0,0) |- node[pos=0.34, left, fill=white, inner sep=1pt]{window ACK progress} (schunk.north east);

\draw[flow2] (rbuf.south west) -- ++(-1.4,0) |- node[pos=0.06, xshift=-98mm, above, align=center, fill=white, inner sep=1pt]{buffer chunks then\\in-order reassembly} (spkt.south east);

\node[align=center] at (8.3,-4.05) {Version 10 increases per-packet payload using fixed-size chunks, reducing header overhead\\Receiver buffering prepares transition to selective-repeat style delivery};
\end{tikzpicture}%
}
\caption{Version 10 architecture: chunk-based payload transfer with in-window buffering and ordered reassembly.
}
\label{fig:v10-architecture}
\end{figure}

\clearpage
\section{Version Family D: Selective Repeat and Adaptive Control (v11--v14)}
Family D contains the most advanced protocol behavior in this project.

\subsection{Version 11: Selective-Repeat-Like Buffering}
Receiver stores out-of-order chunks in buffers and delivers data in order when
missing sequence arrives. This avoids unnecessary retransmission of already
received data and improves behavior under reordering.

Key receiver structures include:
\begin{enumerate}
    \item packet-received flags,
    \item chunk data buffer,
    \item chunk length metadata,
    \item expected sequence tracker.
\end{enumerate}

\vspace{1.5\baselineskip}
\begin{figure}[h!]
\centering
\resizebox{\textwidth}{!}{%
\begin{tikzpicture}[
    font=\normalsize,
    sender/.style={draw, rounded corners=3pt, thick, minimum width=5.9cm, minimum height=1.2cm, align=center, fill=blue!10, draw=blue!60!black},
    receiver/.style={draw, rounded corners=3pt, thick, minimum width=5.9cm, minimum height=1.2cm, align=center, fill=teal!10, draw=teal!60!black},
    channel/.style={draw, rounded corners=3pt, thick, minimum width=4.5cm, minimum height=1.05cm, align=center, fill=orange!10, draw=orange!70!black},
    ctrl/.style={-{Latex[length=3.5mm]}, line width=1.0pt, draw=black!80},
    data/.style={-{Latex[length=3.5mm]}, line width=1.25pt, draw=red!70!black},
    ack/.style={-{Latex[length=3.5mm]}, line width=1.25pt, draw=green!55!black},
    buf/.style={-{Latex[length=3.5mm]}, line width=1.25pt, draw=purple!70!black}
]
\node[sender] (swnd) at (-2.4,2.1) {Sender Window State\\(base, next sequence, outstanding set)};
\node[sender] (ssend) at (-2.4,0.4) {Selective Retransmit Sender\\(retransmit only missing sequences)};
\node[sender] (sout) at (-2.4,-1.3) {Packet Output Stream\\(in-order and out-of-order arrivals possible)};

\node[channel] (net) at (8.0,-1.3) {Unreliable Network with Reordering};

\node[receiver] (rin) at (18.4,-1.3) {Receiver Packet Intake\\(seq check against expected tracker)};
\node[receiver] (rbuf) at (18.4,0.4) {Out-of-Order Buffer\\(flags, chunk data, length metadata)};
\node[receiver] (rdel) at (18.4,2.1) {In-Order Delivery\\(release contiguous chunks only)};

\draw[ctrl] (swnd) -- (ssend);
\draw[ctrl] (ssend) -- (sout);
\draw[ctrl] (rin) -- (rbuf);
\draw[ctrl] (rbuf) -- (rdel);

\draw[data] (sout.east) -- node[pos=0.5, above, align=center, fill=white, inner sep=1pt]{DATA packets\\possibly out of order} (net.west);
\draw[data] (net.east) -- node[pos=0.5, above, align=center, fill=white, inner sep=1pt]{arrivals with gaps\\and reordering} (rin.west);

\draw[ack] (rbuf.north west) -- ++(-4.4,0) |- node[pos=0.33, left, align=center, fill=white, inner sep=1pt]{ACK / selective hint\\for missing sequence} (swnd.north east);

\draw[buf] (rbuf.south west) -- ++(-1.6,0) |- node[pos=0.10, xshift=-92mm, above, align=center, fill=white, inner sep=1pt]{buffer out-of-order chunks\\deliver when gap is filled} (ssend.south east);

\node[align=center] at (8.0,-3.45) {Version 11 stores out-of-order chunks and avoids unnecessary retransmission\\Receiver-side tracking enables selective-repeat-like reliability behavior};
\end{tikzpicture}%
}
\caption{Version 11 architecture: selective-repeat-like buffering with out-of-order storage and in-order delivery.
}
\label{fig:v11-architecture}
\end{figure}

\clearpage
\subsection{Version 12: Adaptive RTO Estimation}
Version 12 introduces RTT sampling and adaptive RTO calculation. Fixed timeout
works poorly in variable-delay conditions. Adaptive timeout adjusts retransmit
threshold dynamically and reduces both premature retransmits and long idle
waits.

\vspace{1.5\baselineskip}
\begin{figure}[h!]
\centering
\resizebox{\textwidth}{!}{%
\begin{tikzpicture}[
    font=\normalsize,
    sender/.style={draw, rounded corners=3pt, thick, minimum width=5.9cm, minimum height=1.25cm, align=center, fill=blue!10, draw=blue!60!black},
    receiver/.style={draw, rounded corners=3pt, thick, minimum width=5.9cm, minimum height=1.25cm, align=center, fill=teal!10, draw=teal!60!black},
    channel/.style={draw, rounded corners=3pt, thick, minimum width=4.7cm, minimum height=1.1cm, align=center, fill=orange!10, draw=orange!70!black},
    ctrl/.style={-{Latex[length=3.5mm]}, line width=1.0pt, draw=black!80},
    data/.style={-{Latex[length=3.5mm]}, line width=1.25pt, draw=red!70!black},
    ack/.style={-{Latex[length=3.5mm]}, line width=1.25pt, draw=green!55!black},
    rto/.style={-{Latex[length=3.5mm]}, line width=1.25pt, draw=purple!70!black}
]
\node[sender] (srtt) at (-2.4,2.8) {RTT Sampler and Estimator\\(sampleRTT, estimatedRTT, deviation)};
\node[sender] (srto) at (-2.4,0.5) {Adaptive RTO Controller\\(RTO updated dynamically per ACK)};
\node[sender] (stx) at (-2.4,-1.8) {Sender Transmission Engine\\(start timer, send, retry on timeout)};

\node[channel] (net) at (8.0,-1.8) {Variable-Delay Network Channel};

\node[receiver] (rrx) at (18.4,-1.8) {Receiver Packet Intake\\(validate packet and sequence)};
\node[receiver] (rack) at (18.4,0.5) {ACK Generator with Timing Signal\\(return ACK with receive timing)};
\node[receiver] (rapp) at (18.4,2.8) {Delivery Layer\\(stable throughput under delay variation)};

\draw[ctrl] (srtt) -- (srto);
\draw[ctrl] (srto) -- (stx);
\draw[ctrl] (rrx) -- (rack);
\draw[ctrl] (rack) -- (rapp);

\draw[data] (stx.east) -- node[pos=0.5, above, align=center, fill=white, inner sep=1pt]{\small DATA with adaptive timer\\instead of fixed timeout} (net.west);
\draw[data] (net.east) -- node[pos=0.5, above, align=center, fill=white, inner sep=1pt]{variable-delay delivery\\(jitter and queueing)} (rrx.west);

\draw[ack] (rack.north west) -- ++(-4.6,0) |- node[pos=0.32, left, align=center, fill=white, inner sep=1pt]{ACK plus RTT feedback\\for timer update} (srtt.north east);

\draw[rto] (srto.south east) -- ++(2.0,0) |- node[pos=0.46, below, yshift=-2mm, align=center, fill=white, inner sep=1pt]{adaptive RTO update\\avoid early and late retries} (rrx.south west);

\node[align=center] at (8.0,-4.45) {Version 12 adapts timeout from observed RTT behavior under changing delay\\This reduces premature retransmits and long idle waiting periods};
\end{tikzpicture}%
}
\caption{Version 12 architecture: RTT sampling with adaptive RTO control for variable-delay conditions.
}
\label{fig:v12-architecture}
\end{figure}

\clearpage
\subsection{Version 13: Congestion Window and Slow Start}
This level introduces congestion control variables:
\begin{enumerate}
    \item \texttt{cwnd} (congestion window),
    \item \texttt{ssthresh} (slow start threshold).
\end{enumerate}

Behavior follows a simplified policy:
\begin{enumerate}
    \item grow quickly in slow start,
    \item grow gradually in congestion avoidance,
    \item reduce sending rate after loss indication.
\end{enumerate}

This improves fairness and stability compared to aggressive fixed-window
sending.

\vspace{1.5\baselineskip}
\begin{figure}[h!]
\centering
\resizebox{\textwidth}{!}{%
\begin{tikzpicture}[
    font=\normalsize,
    sender/.style={draw, rounded corners=3pt, thick, minimum width=6.2cm, minimum height=1.35cm, align=center, fill=blue!10, draw=blue!60!black},
    receiver/.style={draw, rounded corners=3pt, thick, minimum width=6.2cm, minimum height=1.35cm, align=center, fill=teal!10, draw=teal!60!black},
    channel/.style={draw, rounded corners=3pt, thick, minimum width=5.0cm, minimum height=1.2cm, align=center, fill=orange!10, draw=orange!70!black},
    ctrl/.style={-{Latex[length=3.5mm]}, line width=1.0pt, draw=black!80},
    data/.style={-{Latex[length=3.7mm]}, line width=1.35pt, draw=red!70!black},
    ack/.style={-{Latex[length=3.7mm]}, line width=1.35pt, draw=green!55!black},
    cong/.style={-{Latex[length=3.7mm]}, line width=1.35pt, draw=purple!70!black}
]
\node[sender] (scwnd) at (-3.2,3.0) {Congestion Controller\\(cwnd and ssthresh state)};
\node[sender] (sgrow) at (-3.2,0.2) {Growth Policy Engine\\(slow start then congestion avoidance)};
\node[sender] (stx) at (-3.2,-2.6) {Transmission Pacer\\(send rate derived from cwnd)};

\node[channel] (net) at (8.8,-2.6) {Bottleneck-Affected Network Path};

\node[receiver] (rrx) at (20.8,-2.6) {Receiver Intake\\(normal ACK stream under healthy flow)};
\node[receiver] (rack) at (20.8,0.2) {ACK Feedback Source\\(acknowledgement clocking)};
\node[receiver] (robs) at (20.8,3.0) {Loss and Delay Observation\\(congestion signal emergence)};

\draw[ctrl] (scwnd) -- (sgrow);
\draw[ctrl] (sgrow) -- (stx);
\draw[ctrl] (rrx) -- (rack);
\draw[ctrl] (rack) -- (robs);

\draw[data] (stx.east) -- node[pos=0.5, above, align=center, fill=white, inner sep=1pt]{window-controlled DATA flow\\rate scales with cwnd} (net.west);
\draw[data] (net.east) -- node[pos=0.5, above, align=center, fill=white, inner sep=1pt]{throughput to receiver\\under queue and delay pressure} (rrx.west);

\draw[ack] (rack.north west) -- ++(-6.4,0) |- node[pos=0.34, left, align=center, fill=white, inner sep=1pt]{ACK clocking updates\\growth phase and pacing} (scwnd.north east);

\draw[cong] (robs.south west) -- ++(-2.8,0) |- node[pos=0.48, above, align=center, fill=white, inner sep=1pt]{loss indication cut rate\\adjust ssthresh and cwnd} (sgrow.south east);

\node[align=center] at (8.8,-5.35) {Version 13 introduces congestion-aware growth and backoff using cwnd and ssthresh\\Fast growth in slow start transitions to stable additive growth under congestion avoidance};
\end{tikzpicture}%
}
\caption{Version 13 architecture: congestion window growth, ACK clocking, and loss-triggered rate reduction.
}
\label{fig:v13-architecture}
\end{figure}

\clearpage
\subsection{Version 14: Fast Retransmit and Fast Recovery}
Version 14 is the most complete implementation within Family D.
The project then continues through versions 15--24, which extend this baseline
with more advanced adaptive and Reno-like refinements.
Enhancements include:

\begin{enumerate}
    \item duplicate ACK counting,
    \item triple-duplicate-ACK loss inference,
    \item fast retransmission without waiting for timeout,
    \item temporary fast-recovery window inflation and controlled deflation.
\end{enumerate}

This behavior closely resembles a learning-scale version of TCP Reno.

\vspace{1.5\baselineskip}
\begin{figure}[h!]
\centering
\resizebox{\textwidth}{!}{%
\begin{tikzpicture}[
    font=\normalsize,
    sender/.style={draw, rounded corners=3pt, thick, minimum width=6.0cm, minimum height=1.2cm, align=center, fill=blue!10, draw=blue!60!black},
    receiver/.style={draw, rounded corners=3pt, thick, minimum width=6.0cm, minimum height=1.2cm, align=center, fill=teal!10, draw=teal!60!black},
    channel/.style={draw, rounded corners=3pt, thick, minimum width=4.8cm, minimum height=1.1cm, align=center, fill=orange!10, draw=orange!70!black},
    ctrl/.style={-{Latex[length=3.6mm]}, line width=1.0pt, draw=black!80},
    data/.style={-{Latex[length=3.8mm]}, line width=1.35pt, draw=red!70!black},
    ack/.style={-{Latex[length=3.8mm]}, line width=1.35pt, draw=green!55!black},
    fast/.style={-{Latex[length=3.8mm]}, line width=1.35pt, draw=purple!70!black}
]
\node[sender] (sdup) at (-4.2,3.6) {Duplicate ACK Tracker\\(count duplicate ACKs and trigger threshold)};
\node[sender] (sfr) at (-4.2,0.4) {Fast Retransmit and Recovery Logic\\(triple-duplicate ACK -> immediate retransmit)};
\node[sender] (stx) at (-4.2,-2.8) {Sender Transmission Engine\\(recovery window inflate then controlled deflate)};

\node[channel] (net) at (9.2,-2.8) {Loss-Prone Network Path};

\node[receiver] (rrx) at (22.2,-2.8) {Receiver Intake\\(in-order and out-of-order packet observation)};
\node[receiver] (rack) at (22.2,0.4) {ACK Stream Generator\\(duplicate ACKs on missing sequence)};
\node[receiver] (robs) at (22.2,3.6) {Loss Signal Source\\(duplicate ACK pattern visible to sender)};

\draw[ctrl] (sdup) -- (sfr);
\draw[ctrl] (sfr) -- (stx);
\draw[ctrl] (rrx) -- (rack);
\draw[ctrl] (rack) -- (robs);

\draw[data] (stx.east) -- node[pos=0.5, above, align=center, fill=white, inner sep=1pt]{DATA flow with recovery-aware pacing\\avoid waiting for timeout on loss} (net.west);
\draw[data] (net.east) -- node[pos=0.5, above, align=center, fill=white, inner sep=1pt]{delivery stream with occasional loss\\receiver observes missing sequence} (rrx.west);

\draw[ack] (rack.north west) -- ++(-7.2,0) |- node[pos=0.34, left, align=center, fill=white, inner sep=1pt]{duplicate ACK stream\\drives sender loss inference} (sdup.north east);

\draw[fast] (robs.south west) -- ++(-3.2,0) |- node[pos=0.50, above, align=center, fill=white, inner sep=1pt]{triple-duplicate ACK to fast retransmit\\fast recovery inflate then deflate cwnd} (sfr.south east);

\node[align=center] at (9.2,-6.1) {Version 14 adds fast retransmit and fast recovery to react before timeout\\This improves responsiveness while maintaining controlled congestion behavior};
\end{tikzpicture}%
}
\caption{Version 14 architecture: duplicate-ACK loss inference with fast retransmit and fast recovery control.
}
\label{fig:v14-architecture}
\end{figure}

\clearpage
\section{Version Family E: Adaptive and Reno-Like Extensions (v15--v24)}
The later versions extend the same transport story into a more advanced and
educational TCP-like space. This family focuses on RTT-sensitive timeout
refinement, timestamp-oriented state awareness, SACK-like reasoning,
congestion-avoidance behavior, and improved recovery handling. The visualizer
also uses these later versions to show how the FSM and UML timeline can be used
to explain protocol decisions step by step.

The primary design inspiration for this family comes from the RFCs and books
that describe TCP evolution in practice. In particular, RFC 6298 explains the
retransmission timer, RFC 5681 covers congestion control behavior, RFC 2018
describes selective acknowledgment options, RFC 7323 covers TCP extensions for
high performance, and RFC 793 provides the core TCP model. The implementation
style is also informed by Stevens, Kurose and Ross, and Tanenbaum and
Wetherall, because those texts help connect protocol theory with the concrete
sender and receiver code structure used in the project \cite{rfc6298,rfc5681,rfc2018,rfc7323,rfc793,stevens_tcpip,kurose_ross,tanenbaum_wetherall}.

This family is especially valuable in the report because it shows the point at
which the learning project becomes visually and conceptually close to real TCP
thinking, while still remaining simple enough for students to trace manually.

\clearpage
\subsection{Version 15: RTT-Smoothed Timeout Refinement}
Version 15 stabilizes retransmission timing using smoothed RTT and deviation
tracking. Instead of relying on coarse timeout behavior, this version updates
RTO with a more disciplined estimator so sender retries are less noisy.

\begin{figure}[h!]
\centering
\resizebox{\textwidth}{!}{%
\begin{tikzpicture}[
    font=\normalsize,
    sender/.style={draw, rounded corners=3pt, thick, minimum width=6.2cm, minimum height=1.25cm, align=center, fill=blue!10, draw=blue!60!black},
    receiver/.style={draw, rounded corners=3pt, thick, minimum width=6.2cm, minimum height=1.25cm, align=center, fill=teal!10, draw=teal!60!black},
    channel/.style={draw, rounded corners=3pt, thick, minimum width=4.8cm, minimum height=1.1cm, align=center, fill=orange!10, draw=orange!70!black},
    data/.style={-{Latex[length=3.5mm]}, line width=1.2pt, draw=red!70!black},
    ack/.style={-{Latex[length=3.5mm]}, line width=1.2pt, draw=green!55!black},
    ctrl/.style={-{Latex[length=3.5mm]}, line width=1.1pt, draw=purple!70!black}
]
\node[sender] (s1) at (-3.0,2.1) {RTT Sampler\\(sampleRTT and ACK timing)};
\node[sender] (s2) at (-3.0,-0.2) {RTO Updater\\(smoothed timer calculation)};
\node[sender] (s3) at (-3.0,-2.5) {Send/Retry Engine\\(retry threshold from adaptive RTO)};

\node[channel] (net) at (8.0,-2.5) {Variable-Delay Path};

\node[receiver] (r1) at (19.0,-2.5) {Receiver Delivery Pipeline};
\node[receiver] (r2) at (19.0,-0.2) {ACK Timing Feedback};
\node[receiver] (r3) at (19.0,2.1) {Stable ACK Clock Source};

\draw[ctrl] (s1) -- (s2);
\draw[ctrl] (s2) -- (s3);
\draw[ctrl] (r1) -- (r2);
\draw[ctrl] (r2) -- (r3);

\draw[data] (s3.east) -- node[pos=0.5, above, fill=white, inner sep=1pt]{timed DATA transmission} (net.west);
\draw[data] (net.east) -- node[pos=0.5, above, fill=white, inner sep=1pt]{delay-affected arrival stream} (r1.west);
\draw[ack] (r2.north west) -- ++(-5.0,0) |- node[pos=0.34, left, fill=white, inner sep=1pt]{ACK timing} (s1.north east);
\end{tikzpicture}%
}
\caption{Version 15 architecture: smoothed RTT-driven retransmission timeout refinement.}
\label{fig:v15-architecture}
\end{figure}

\clearpage
\subsection{Version 16: Adaptive Window Discipline}
Version 16 connects timeout behavior with sender pacing and window discipline.
When timing uncertainty grows, send aggressiveness is reduced. When ACK rhythm
stabilizes, throughput is allowed to grow safely.

\begin{figure}[h!]
\centering
\resizebox{\textwidth}{!}{%
\begin{tikzpicture}[
    font=\normalsize,
    sender/.style={draw, rounded corners=3pt, thick, minimum width=6.0cm, minimum height=1.2cm, align=center, fill=blue!10, draw=blue!60!black},
    receiver/.style={draw, rounded corners=3pt, thick, minimum width=6.0cm, minimum height=1.2cm, align=center, fill=teal!10, draw=teal!60!black},
    channel/.style={draw, rounded corners=3pt, thick, minimum width=4.8cm, minimum height=1.05cm, align=center, fill=orange!10, draw=orange!70!black},
    data/.style={-{Latex[length=3.5mm]}, line width=1.2pt, draw=red!70!black},
    ack/.style={-{Latex[length=3.5mm]}, line width=1.2pt, draw=green!55!black},
    flow/.style={-{Latex[length=3.5mm]}, line width=1.1pt, draw=purple!70!black}
]
\node[sender] (s1) at (-3.2,2.0) {Adaptive Sender Window\\(base, next, pacing gate)};
\node[sender] (s2) at (-3.2,-0.3) {Timeout-Aware Rate Control};
\node[sender] (s3) at (-3.2,-2.6) {Windowed Data Output};

\node[channel] (net) at (8.2,-2.6) {Loss and Delay Channel};

\node[receiver] (r1) at (19.6,-2.6) {Packet Intake and Ordering};
\node[receiver] (r2) at (19.6,-0.3) {ACK and Expected-Seq State};
\node[receiver] (r3) at (19.6,2.0) {Flow Stability Feedback};

\draw[flow] (s1) -- (s2);
\draw[flow] (s2) -- (s3);
\draw[flow] (r1) -- (r2);
\draw[flow] (r2) -- (r3);

\draw[data] (s3.east) -- node[pos=0.5, above, fill=white, inner sep=1pt]{adaptive window DATA} (net.west);
\draw[data] (net.east) -- node[pos=0.5, above, fill=white, inner sep=1pt]{receiver in-order processing} (r1.west);
\draw[ack] (r2.north west) -- ++(-5.6,0) |- node[pos=0.36, left, fill=white, inner sep=1pt]{ACK rhythm and gap hints} (s1.north east);
\end{tikzpicture}%
}
\caption{Version 16 architecture: timeout-aware adaptive window control.}
\label{fig:v16-architecture}
\end{figure}

\clearpage
\subsection{Version 17: Timestamp Option Integration}
Version 17 introduces timestamp-style metadata in packet exchange so sender can
infer timing more directly and avoid ambiguous retransmission decisions.

\begin{figure}[h!]
\centering
\resizebox{\textwidth}{!}{%
\begin{tikzpicture}[
    font=\normalsize,
    sender/.style={draw, rounded corners=3pt, thick, minimum width=6.0cm, minimum height=1.2cm, align=center, fill=blue!10, draw=blue!60!black},
    receiver/.style={draw, rounded corners=3pt, thick, minimum width=6.0cm, minimum height=1.2cm, align=center, fill=teal!10, draw=teal!60!black},
    channel/.style={draw, rounded corners=3pt, thick, minimum width=4.7cm, minimum height=1.05cm, align=center, fill=orange!10, draw=orange!70!black},
    data/.style={-{Latex[length=3.5mm]}, line width=1.2pt, draw=red!70!black},
    ack/.style={-{Latex[length=3.5mm]}, line width=1.2pt, draw=green!55!black},
    opt/.style={-{Latex[length=3.5mm]}, line width=1.1pt, draw=purple!70!black}
]
\node[sender] (s1) at (-3.0,1.8) {Timestamp Inserter\\(ts\_val in outgoing packet)};
\node[sender] (s2) at (-3.0,-0.4) {RTT Estimator with TS Echo};
\node[sender] (s3) at (-3.0,-2.6) {Retransmit Decision Logic};

\node[channel] (net) at (8.0,-2.6) {Transport Path};

\node[receiver] (r1) at (19.0,-2.6) {Timestamp Parser};
\node[receiver] (r2) at (19.0,-0.4) {ACK Builder with ts\_ecr};
\node[receiver] (r3) at (19.0,1.8) {Option-Aware Receiver State};

\draw[opt] (s1) -- (s2);
\draw[opt] (s2) -- (s3);
\draw[opt] (r1) -- (r2);
\draw[opt] (r2) -- (r3);

\draw[data] (s3.east) -- node[pos=0.5, above, fill=white, inner sep=1pt]{DATA + timestamp option} (net.west);
\draw[data] (net.east) -- node[pos=0.5, above, fill=white, inner sep=1pt]{option-aware reception} (r1.west);
\draw[ack] (r2.north west) -- ++(-5.0,0) |- node[pos=0.35, left, fill=white, inner sep=1pt]{ACK with ts echo} (s1.north east);
\end{tikzpicture}%
}
\caption{Version 17 architecture: timestamp insertion and echo-aware feedback loop.}
\label{fig:v17-architecture}
\end{figure}

\clearpage
\subsection{Version 18: SACK-Inspired Receiver Reporting}
Version 18 introduces selective acknowledgment style hints so sender can see
which blocks arrived beyond a gap. This reduces unnecessary retransmission of
already received data.

\begin{figure}[h!]
\centering
\resizebox{\textwidth}{!}{%
\begin{tikzpicture}[
    font=\normalsize,
    sender/.style={draw, rounded corners=3pt, thick, minimum width=6.0cm, minimum height=1.2cm, align=center, fill=blue!10, draw=blue!60!black},
    receiver/.style={draw, rounded corners=3pt, thick, minimum width=6.0cm, minimum height=1.2cm, align=center, fill=teal!10, draw=teal!60!black},
    channel/.style={draw, rounded corners=3pt, thick, minimum width=4.8cm, minimum height=1.05cm, align=center, fill=orange!10, draw=orange!70!black},
    data/.style={-{Latex[length=3.5mm]}, line width=1.2pt, draw=red!70!black},
    ack/.style={-{Latex[length=3.5mm]}, line width=1.2pt, draw=green!55!black},
    sack/.style={-{Latex[length=3.5mm]}, line width=1.1pt, draw=purple!70!black}
]
\node[sender] (s1) at (-3.0,2.0) {Gap Map Tracker\\(missing sequence focus)};
\node[sender] (s2) at (-3.0,-0.3) {Selective Retransmit Planner};
\node[sender] (s3) at (-3.0,-2.6) {Loss-Targeted Packet Output};

\node[channel] (net) at (8.0,-2.6) {Reordering and Loss Channel};

\node[receiver] (r1) at (19.0,-2.6) {Out-of-Order Buffer State};
\node[receiver] (r2) at (19.0,-0.3) {SACK Block Builder};
\node[receiver] (r3) at (19.0,2.0) {ACK + Block Range Feedback};

\draw[sack] (s1) -- (s2);
\draw[sack] (s2) -- (s3);
\draw[sack] (r1) -- (r2);
\draw[sack] (r2) -- (r3);

\draw[data] (s3.east) -- node[pos=0.5, above, fill=white, inner sep=1pt]{windowed DATA stream} (net.west);
\draw[data] (net.east) -- node[pos=0.5, above, fill=white, inner sep=1pt]{buffered arrivals with gaps} (r1.west);
\draw[ack] (r3.west) -- ++(-5.2,0) |- node[pos=0.36, left, fill=white, inner sep=1pt]{ACK + SACK-like block hints} (s1.north east);
\end{tikzpicture}%
}
\caption{Version 18 architecture: SACK-inspired gap reporting and selective resend planning.}
\label{fig:v18-architecture}
\end{figure}

\clearpage
\subsection{Version 19: Congestion-Avoidance Stabilization}
Version 19 strengthens congestion-avoidance transitions. The sender responds to
loss and delay pressure with smoother cwnd behavior, reducing oscillation while
keeping reasonable throughput.

\begin{figure}[h!]
\centering
\resizebox{\textwidth}{!}{%
\begin{tikzpicture}[
    font=\normalsize,
    sender/.style={draw, rounded corners=3pt, thick, minimum width=6.1cm, minimum height=1.2cm, align=center, fill=blue!10, draw=blue!60!black},
    receiver/.style={draw, rounded corners=3pt, thick, minimum width=6.1cm, minimum height=1.2cm, align=center, fill=teal!10, draw=teal!60!black},
    channel/.style={draw, rounded corners=3pt, thick, minimum width=4.8cm, minimum height=1.05cm, align=center, fill=orange!10, draw=orange!70!black},
    data/.style={-{Latex[length=3.5mm]}, line width=1.2pt, draw=red!70!black},
    ack/.style={-{Latex[length=3.5mm]}, line width=1.2pt, draw=green!55!black},
    cong/.style={-{Latex[length=3.5mm]}, line width=1.1pt, draw=purple!70!black}
]
\node[sender] (s1) at (-3.2,2.0) {cwnd/ssthresh Controller};
\node[sender] (s2) at (-3.2,-0.3) {Slow-Start to CA Transition Logic};
\node[sender] (s3) at (-3.2,-2.6) {Congestion-Aware Send Pacer};

\node[channel] (net) at (8.2,-2.6) {Bottleneck-Sensitive Link};

\node[receiver] (r1) at (19.6,-2.6) {Data Receiver Pipeline};
\node[receiver] (r2) at (19.6,-0.3) {ACK Clock and DupACK Signal};
\node[receiver] (r3) at (19.6,2.0) {Congestion Signal Visibility};

\draw[cong] (s1) -- (s2);
\draw[cong] (s2) -- (s3);
\draw[cong] (r1) -- (r2);
\draw[cong] (r2) -- (r3);

\draw[data] (s3.east) -- node[pos=0.5, above, fill=white, inner sep=1pt]{rate-shaped DATA} (net.west);
\draw[data] (net.east) -- node[pos=0.5, above, fill=white, inner sep=1pt]{queue-influenced delivery} (r1.west);
\draw[ack] (r2.north west) -- ++(-5.8,0) |- node[pos=0.35, left, fill=white, inner sep=1pt]{ACK feedback and DupACK events} (s1.north east);
\end{tikzpicture}%
}
\caption{Version 19 architecture: stabilized congestion-avoidance behavior with ACK-clocked feedback.}
\label{fig:v19-architecture}
\end{figure}

\clearpage
\subsection{Version 20: Loss-Recovery Coupled with Option Feedback}
Version 20 combines fast recovery behavior with richer feedback from timestamp
and selective-ack style information. Recovery targets missing data more
precisely and returns to normal transmission with better control.

\begin{figure}[h!]
\centering
\resizebox{\textwidth}{!}{%
\begin{tikzpicture}[
    font=\normalsize,
    sender/.style={draw, rounded corners=3pt, thick, minimum width=6.2cm, minimum height=1.2cm, align=center, fill=blue!10, draw=blue!60!black},
    receiver/.style={draw, rounded corners=3pt, thick, minimum width=6.2cm, minimum height=1.2cm, align=center, fill=teal!10, draw=teal!60!black},
    channel/.style={draw, rounded corners=3pt, thick, minimum width=4.9cm, minimum height=1.05cm, align=center, fill=orange!10, draw=orange!70!black},
    data/.style={-{Latex[length=3.5mm]}, line width=1.2pt, draw=red!70!black},
    ack/.style={-{Latex[length=3.5mm]}, line width=1.2pt, draw=green!55!black},
    rec/.style={-{Latex[length=3.5mm]}, line width=1.1pt, draw=purple!70!black}
]
\node[sender] (s1) at (-3.3,2.0) {Recovery Trigger Logic\\(DupACK and option hints)};
\node[sender] (s2) at (-3.3,-0.3) {Selective Fast-Retransmit Unit};
\node[sender] (s3) at (-3.3,-2.6) {Controlled Exit to Normal Flow};

\node[channel] (net) at (8.3,-2.6) {Mixed Loss/Reorder Path};

\node[receiver] (r1) at (19.9,-2.6) {Gap-Aware Receive Buffer};
\node[receiver] (r2) at (19.9,-0.3) {ACK + TS + Block Feedback};
\node[receiver] (r3) at (19.9,2.0) {Recovery Guidance Signal};

\draw[rec] (s1) -- (s2);
\draw[rec] (s2) -- (s3);
\draw[rec] (r1) -- (r2);
\draw[rec] (r2) -- (r3);

\draw[data] (s3.east) -- node[pos=0.5, above, fill=white, inner sep=1pt]{recovery-phase DATA stream} (net.west);
\draw[data] (net.east) -- node[pos=0.5, above, fill=white, inner sep=1pt]{gap-observing receiver path} (r1.west);
\draw[ack] (r2.north west) -- ++(-5.9,0) |- node[pos=0.34, left, fill=white, inner sep=1pt]{rich ACK feedback} (s1.north east);
\end{tikzpicture}%
}
\caption{Version 20 architecture: fast recovery guided by richer feedback signals.}
\label{fig:v20-architecture}
\end{figure}

\clearpage
\subsection{Version 21: Sender State Consolidation}
Version 21 consolidates sender logic for timers, congestion state, option
tracking, and retransmission queues so control decisions remain consistent
across difficult traffic patterns.

\begin{figure}[h!]
\centering
\resizebox{\textwidth}{!}{%
\begin{tikzpicture}[
    font=\normalsize,
    sender/.style={draw, rounded corners=3pt, thick, minimum width=6.1cm, minimum height=1.2cm, align=center, fill=blue!10, draw=blue!60!black},
    receiver/.style={draw, rounded corners=3pt, thick, minimum width=6.1cm, minimum height=1.2cm, align=center, fill=teal!10, draw=teal!60!black},
    channel/.style={draw, rounded corners=3pt, thick, minimum width=4.8cm, minimum height=1.05cm, align=center, fill=orange!10, draw=orange!70!black},
    data/.style={-{Latex[length=3.5mm]}, line width=1.2pt, draw=red!70!black},
    ack/.style={-{Latex[length=3.5mm]}, line width=1.2pt, draw=green!55!black},
    ctrl/.style={-{Latex[length=3.5mm]}, line width=1.1pt, draw=purple!70!black}
]
\node[sender] (s1) at (-3.2,2.0) {Unified Sender Control State};
\node[sender] (s2) at (-3.2,-0.3) {Timer + Congestion + Option Arbiter};
\node[sender] (s3) at (-3.2,-2.6) {Consistent Tx/Rtx Queue Manager};

\node[channel] (net) at (8.2,-2.6) {Transport Channel};

\node[receiver] (r1) at (19.6,-2.6) {Ordered Delivery Core};
\node[receiver] (r2) at (19.6,-0.3) {ACK Signal Generator};
\node[receiver] (r3) at (19.6,2.0) {State Feedback Source};

\draw[ctrl] (s1) -- (s2);
\draw[ctrl] (s2) -- (s3);
\draw[ctrl] (r1) -- (r2);
\draw[ctrl] (r2) -- (r3);

\draw[data] (s3.east) -- node[pos=0.5, above, fill=white, inner sep=1pt]{state-consistent DATA stream} (net.west);
\draw[data] (net.east) -- node[pos=0.5, above, fill=white, inner sep=1pt]{receiver acceptance pipeline} (r1.west);
\draw[ack] (r2.north west) -- ++(-5.7,0) |- node[pos=0.34, left, fill=white, inner sep=1pt]{ACK and state feedback} (s1.north east);
\end{tikzpicture}%
}
\caption{Version 21 architecture: consolidated sender control for stable decision flow.}
\label{fig:v21-architecture}
\end{figure}

\clearpage
\subsection{Version 22: Receiver Reassembly and Delivery Hardening}
Version 22 strengthens the receiver side by improving buffering, gap filling,
and reassembly commit rules. This reduces edge-case delivery mistakes when
packets arrive late, duplicated, or out of order.

\begin{figure}[h!]
\centering
\resizebox{\textwidth}{!}{%
\begin{tikzpicture}[
    font=\normalsize,
    sender/.style={draw, rounded corners=3pt, thick, minimum width=6.0cm, minimum height=1.2cm, align=center, fill=blue!10, draw=blue!60!black},
    receiver/.style={draw, rounded corners=3pt, thick, minimum width=6.0cm, minimum height=1.2cm, align=center, fill=teal!10, draw=teal!60!black},
    channel/.style={draw, rounded corners=3pt, thick, minimum width=4.8cm, minimum height=1.05cm, align=center, fill=orange!10, draw=orange!70!black},
    data/.style={-{Latex[length=3.5mm]}, line width=1.2pt, draw=red!70!black},
    ack/.style={-{Latex[length=3.5mm]}, line width=1.2pt, draw=green!55!black},
    rec/.style={-{Latex[length=3.5mm]}, line width=1.1pt, draw=purple!70!black}
]
\node[sender] (s1) at (-3.0,2.0) {Sender Window Output};
\node[sender] (s2) at (-3.0,-0.3) {Selective Resend Support};
\node[sender] (s3) at (-3.0,-2.6) {Delivery-Oriented Packet Stream};

\node[channel] (net) at (8.0,-2.6) {Uncertain Arrival Channel};

\node[receiver] (r1) at (19.0,-2.6) {Reassembly Buffer Matrix};
\node[receiver] (r2) at (19.0,-0.3) {Gap Fill and Contiguous Commit Logic};
\node[receiver] (r3) at (19.0,2.0) {Reliable App Delivery Stage};

\draw[rec] (s1) -- (s2);
\draw[rec] (s2) -- (s3);
\draw[rec] (r1) -- (r2);
\draw[rec] (r2) -- (r3);

\draw[data] (s3.east) -- node[pos=0.5, above, fill=white, inner sep=1pt]{windowed DATA packets} (net.west);
\draw[data] (net.east) -- node[pos=0.5, above, fill=white, inner sep=1pt]{out-of-order chunk arrivals} (r1.west);
\draw[ack] (r2.north west) -- ++(-5.0,0) |- node[pos=0.34, left, fill=white, inner sep=1pt]{ACK for committed ranges} (s1.north east);
\end{tikzpicture}%
}
\caption{Version 22 architecture: hardened receiver reassembly and contiguous delivery control.}
\label{fig:v22-architecture}
\end{figure}

\clearpage
\subsection{Version 23: End-to-End Robustness Tuning}
Version 23 tunes interaction among timer control, congestion response,
selective retransmission, and receiver commit rules. The focus is robust
behavior under mixed delay, reordering, and burst loss.

\begin{figure}[h!]
\centering
\resizebox{\textwidth}{!}{%
\begin{tikzpicture}[
    font=\normalsize,
    sender/.style={draw, rounded corners=3pt, thick, minimum width=6.2cm, minimum height=1.2cm, align=center, fill=blue!10, draw=blue!60!black},
    receiver/.style={draw, rounded corners=3pt, thick, minimum width=6.2cm, minimum height=1.2cm, align=center, fill=teal!10, draw=teal!60!black},
    channel/.style={draw, rounded corners=3pt, thick, minimum width=4.8cm, minimum height=1.05cm, align=center, fill=orange!10, draw=orange!70!black},
    data/.style={-{Latex[length=3.5mm]}, line width=1.2pt, draw=red!70!black},
    ack/.style={-{Latex[length=3.5mm]}, line width=1.2pt, draw=green!55!black},
    tune/.style={-{Latex[length=3.5mm]}, line width=1.1pt, draw=purple!70!black}
]
\node[sender] (s1) at (-3.2,2.0) {Integrated Robustness Policy};
\node[sender] (s2) at (-3.2,-0.3) {Timer/CWND/Resend Coordinator};
\node[sender] (s3) at (-3.2,-2.6) {Resilient Packet Scheduler};

\node[channel] (net) at (8.2,-2.6) {Burst-Loss and Reorder Channel};

\node[receiver] (r1) at (19.6,-2.6) {Robust Buffer and Delivery Core};
\node[receiver] (r2) at (19.6,-0.3) {Feedback with Gap/Timing Context};
\node[receiver] (r3) at (19.6,2.0) {Application-Safe Commit Stage};

\draw[tune] (s1) -- (s2);
\draw[tune] (s2) -- (s3);
\draw[tune] (r1) -- (r2);
\draw[tune] (r2) -- (r3);

\draw[data] (s3.east) -- node[pos=0.5, above, fill=white, inner sep=1pt]{robustness-tuned DATA flow} (net.west);
\draw[data] (net.east) -- node[pos=0.5, above, fill=white, inner sep=1pt]{mixed-failure delivery path} (r1.west);
\draw[ack] (r2.north west) -- ++(-5.8,0) |- node[pos=0.35, left, fill=white, inner sep=1pt]{context-rich ACK feedback} (s1.north east);
\end{tikzpicture}%
}
\caption{Version 23 architecture: coordinated end-to-end robustness tuning under mixed failures.}
\label{fig:v23-architecture}
\end{figure}

\clearpage
\subsection{Version 24: Final Integrated Transport Stage}
Version 24 is the final integrated stage of the project. It combines the full
set of educational transport features: adaptive timing, congestion-aware
windows, timestamp-aware feedback, SACK-inspired loss hints, and controlled
recovery behavior.

\begin{figure}[h!]
\centering
\resizebox{\textwidth}{!}{%
\begin{tikzpicture}[
    font=\normalsize,
    sender/.style={draw, rounded corners=3pt, thick, minimum width=6.3cm, minimum height=1.25cm, align=center, fill=blue!10, draw=blue!60!black},
    receiver/.style={draw, rounded corners=3pt, thick, minimum width=6.3cm, minimum height=1.25cm, align=center, fill=teal!10, draw=teal!60!black},
    channel/.style={draw, rounded corners=3pt, thick, minimum width=5.0cm, minimum height=1.1cm, align=center, fill=orange!10, draw=orange!70!black},
    data/.style={-{Latex[length=3.7mm]}, line width=1.25pt, draw=red!70!black},
    ack/.style={-{Latex[length=3.7mm]}, line width=1.25pt, draw=green!55!black},
    all/.style={-{Latex[length=3.7mm]}, line width=1.15pt, draw=purple!70!black}
]
\node[sender] (s1) at (-3.4,2.2) {Integrated Sender Brain\\(RTO + cwnd + recovery + option parsing)};
\node[sender] (s2) at (-3.4,-0.2) {Selective and Fast Recovery Engine};
\node[sender] (s3) at (-3.4,-2.6) {Stable High-Utility Packet Output};

\node[channel] (net) at (8.4,-2.6) {Full Experimental Transport Channel};

\node[receiver] (r1) at (20.2,-2.6) {Gap-Tolerant Receive and Reassembly};
\node[receiver] (r2) at (20.2,-0.2) {ACK + TS + SACK-Style Feedback Layer};
\node[receiver] (r3) at (20.2,2.2) {Reliable In-Order Delivery to App Layer};

\draw[all] (s1) -- (s2);
\draw[all] (s2) -- (s3);
\draw[all] (r1) -- (r2);
\draw[all] (r2) -- (r3);

\draw[data] (s3.east) -- node[pos=0.5, above, align=center, fill=white, inner sep=1pt]{integrated DATA flow with\\adaptive transport control} (net.west);
\draw[data] (net.east) -- node[pos=0.5, above, align=center, fill=white, inner sep=1pt]{robust receiver intake with\\ordered reconstruction} (r1.west);
\draw[ack] (r2.north west) -- ++(-6.0,0) |- node[pos=0.34, left, align=center, fill=white, inner sep=1pt]{combined ACK, timing,\\and loss-hint feedback} (s1.north east);
\end{tikzpicture}%
}
\caption{Version 24 architecture: final integrated transport design across reliability, timing, and recovery.}
\label{fig:v24-architecture}
\end{figure}

\clearpage
\subsection{Family E Summary}
Versions 15--24 complete the transition from a staged classroom transport model
to an integrated TCP-inspired system. The sequence shows how timing, selective
feedback, congestion awareness, and structured recovery can be layered without
losing traceability for learners.

\section{Chapter Summary}
This chapter explained how the full project evolves from a basic UDP demo to a
rich transport protocol laboratory. The next chapter formalizes the main
algorithms with equations and analysis so that implementation behavior can be
understood quantitatively.understood quantitatively. Chapter 4 then presents detailed measurements,
comparison tables, and visual timelines to complete the educational narrative.
